\chapter{System Requirement Specification}
\label{ch:system_requirements}

\section{Chapter Overview}
\label{sec:req_chapter_overview}
For the design of a robust, effective, and versatile paper keyboard, a System Requirement Specification has to be established. Visual HCI design without the use of any physical switches should include appropriate systems analysis. The camera tracking, normalization of hand joints, homography estimation, and deep learning touch detection must meet the needs of users \cite{ReviewVirtualKeyboard2020, Maman2023TypeNet}.

This chapter discusses Requirement Specification and UML models for both applications: Application 1 (Layout Designer Suite \texttt{designer/designer\_app.py}) and Application 2 (Virtual Keyboard Runtime Engine and Command Mapper). Stakeholder analysis with user identification, Power-Interest matrix, and design trade-offs is discussed in Section~\ref{sec:req_stakeholder}. Operationalization of the research objectives and test data to derive software requirements is discussed in Section~\ref{sec:req_operationalization}. Comprehensive UML analysis is discussed in Section~\ref{sec:req_uml_analysis}, which consists of Use Case diagrams with formal specifications, Class diagrams with CRC descriptions, Activity diagrams, Sequence diagrams, Statechart diagrams, and Deployment diagrams. System Architecture Framework and Design Patterns are discussed in Section~\ref{sec:req_architecture}. Functional and Non-Functional Requirements are discussed in Section~\ref{sec:req_fn_nfn}. Finally, Section~\ref{sec:req_summary} summarizes the chapter.

\section{Stakeholder Analysis}
\label{sec:req_stakeholder}
Through stakeholder analysis, we will determine who uses the product or stakeholders influenced by the virtual keyboard. Therefore, the system will satisfy the requirements of the stakeholders \cite{ReviewVirtualKeyboard2020}.

\subsection{Stakeholder Identification and Classification}
\label{subsec:req_stakeholder_classification}
The following is a list of four major groups that act as stakeholders in virtual keyboard technology through papers:
\begin{enumerate}
    \item \textbf{Primary End Users:} This group consists of people who type on paper sheets in order to input data or control computer systems. These include regular computer end users, typists, and people who have problems with movement and require large keyboards \cite{Chen2024Gaze}.
    \item \textbf{Subject Specific End Users:} This includes people from particular subjects such as digital audio editors, video editors, industrial machine operators, and CAD designers who require customized macros for software keystrokes \cite{ReviewVirtualKeyboard2020}.
    \item \textbf{Medical and Clinical Professionals:} This category includes doctors, nurses, and lab technicians who work in sterile environments, where normal plastic keyboards can collect dirt and cause germ spreading \cite{ReviewVirtualKeyboard2020}.
    \item \textbf{System Developers and Researchers:} This group includes people conducting research in computer vision, software developers, and system administrators for layout design, XML configuration, neural model training (\texttt{run\_all.py}), and runtime software.
\end{enumerate}

The stakeholders, their responsibilities, and needs have been detailed in Table~\ref{tab:stakeholder_matrix} below.

\begin{table}[htbp]
\centering
\caption{Comprehensive stakeholder identification, operational roles, and key requirements matrix.}
\label{tab:stakeholder_matrix}
\begin{tabular}{l l l l}
\toprule
\textbf{Stakeholder Group} & \textbf{Specific User Category} & \textbf{Operational Role and Context} & \textbf{Key System Requirements} \\
\midrule
\textbf{Primary End-Users} & General Desktop Typists & - Daily text entry and computer navigation & - Low entry latency ($< 30$ ms) \\
 & & - Desktop computing environments & - High touch trigger accuracy ($> 95\%$) \\
 & Accessibility Users & - Motor-impaired or low-dexterity users & - Enlarged key size customization \\
 & & - Customized physical layout sheets & - Low Index of Difficulty ($ID$) under Fitts' Law \\
\midrule
\textbf{Domain Specialists} & Audio / Video Engineers & - Dynamic macro triggering in DAWs & - Multi-key shell script binding \\
 & Industrial Operators & - Machine tool control pads & - Dynamic layout XML hot-swapping \\
\midrule
\textbf{Healthcare / Clinical} & Surgeons and ICU Staff & - Touch interaction in sterile operating rooms & - Disposable, sanitizable paper sheets \\
 & Cleanroom Technicians & - Pathogen-free computer interface & - Zero physical switch contamination \\
\midrule
\textbf{System Engineers} & Vision Researchers & - Algorithm benchmarking (\texttt{run\_all.py}) & - Scale-normalized feature exports \\
 & Software Administrators & - System command mapping and deployment & - Validated XML schema synchronization \\
\bottomrule
\end{tabular}%
\end{table}

\subsection{Stakeholder Power-Interest Matrix}
\label{subsec:req_stakeholder_matrix}
In order to prioritize the requirements, the stakeholders have been mapped in terms of their power and interest in the Power-Interest Matrix as shown in Figure~\ref{fig:stakeholder_power_interest} below:

\begin{figure}[htbp]
\centering

\includegraphics[width=0.85\textwidth]{figures/chapter04_fig01.png}

\caption{Stakeholder Power-Interest Matrix governing project requirement prioritization.}
\label{fig:stakeholder_power_interest}
\end{figure}

As depicted from the above figure (Figure~\ref{fig:stakeholder_power_interest}), System Engineers and Domain Experts are Key Players (High Power and High Interest). On the other hand, Primary Typists and Accessibility Users belong to the High Interest group.

\subsection{Stakeholder Conflict Resolution and Requirement Trade-offs}
\label{subsec:req_conflict_resolution}
Below is the list of problems that emerged because of different interests of various stakeholders and the solutions to those:
\begin{enumerate}
    \item \textbf{Typing Speed vs Computing Power Required:} The person who will be doing typing requires a low-latency device, while the admin requires an application that runs on an ordinary computer with low computations. The solution was found in selecting the video frame rate of 12~FPS (\texttt{resample\_12fps.py}) and using an efficient LSTM network with PyTorch (\texttt{best\_finger\_touch\_lstm.pth}) to make sure that the application is running in real time without requiring any GPU.
    \item \textbf{Custom Layout vs Tracking Accuracy:} Custom macro layouts may have buttons that are placed in the corners of the sheet, thus blocking the tracking markers. The solution was finding the optimal position for the AprilTag markers to be along the outer margins of the sheet, making it possible to find the homography matrix $H$ by using only four corners.
    \item \textbf{Disposal of Papers vs Saving Profiles:} The solution to this problem was saving the geometry of the layout in XML format and making it possible to dispose of or print papers whenever needed without any reconfiguration setup on the computer.
\end{enumerate}

\section{Operationalization Process}
\label{sec:req_operationalization}
The process of operationalization acts as an interface between the high-level goals of the research and the experiments with data collection and low-level software requirements.

\subsection{Mapping Research Data Gathering to System Requirements}
\label{subsec:req_data_mapping}
The system requirements have been determined through the analysis of 57 literature sources, videos collected from 15 users, and measurements of CPU performance.

\begin{table}[htbp]
\centering
\caption{Empirical techniques and data sources used to formulate system requirements.}
\label{tab:empirical_data_requirements}
\begin{tabular}{l l l l}
\toprule
\textbf{Research Objective} & \textbf{Data Gathering / Empirical Technique} & \textbf{Requirement Validity Justification} & \textbf{Mapped System Requirement Module} \\
\midrule
\textbf{Obj 1: Identify Gaps and} & - Systematic literature review (57 papers) & - Establishes baseline latency ($< 30$ ms) & - Core System Architectural Scope \\
Establish Baselines & - Benchmark matrix (\texttt{Table 2.1}) & - Justifies monocular RGB paper approach & - Hardware/Software Constraints \\
\midrule
\textbf{Obj 2: Software Layout} & - User study interviews ($N=15$) & - Confirms need for arbitrary key sizing & - Layout Designer App (\texttt{designer\_app.py}) \\
Designer and PDF/XML Exporter & - Ergonomic Fitts' law evaluation & - Validates synchronized PDF/XML export & - PDF Vector Renderer and XML Exporter \\
\midrule
\textbf{Obj 3: Landmark Pose and} & - 15-participant video dataset ($L_{\text{hand}}$) & - Proves scale invariance across distance $Z$ & - MediaPipe Pose Extractor \\
Scale Normalization Engine & - Distance scaling experiments & - Confirms scale and resolution invariance & - Scale Normalizer (\texttt{normalize\_landmarks.py}) \\
\midrule
\textbf{Obj 4: Deep Temporal} & - Automated benchmark (\texttt{run\_all.py}) & - Benchmarks 22 pure 2D architectures & - Sliding Window Buffer (\texttt{create\_windows.py}) \\
Touch Detection Classifier & - 5-fold cross-validation on 5-frame windows & - Selects optimal PyTorch LSTM model & - PyTorch Classifier (\texttt{best\_finger\_touch\_lstm.pth}) \\
\midrule
\textbf{Obj 5: AprilTag Fiducial} & - Tilt camera experiments ($0^\circ$ to $75^\circ$) & - Solves SVD Planar Homography $H$ & - AprilTag Tracker (\texttt{aprilTag/estimater.py}) \\
Homography Tracking Engine & - Sub-pixel corner localization tests & - Ensures sub-millimeter mapping accuracy & - Planar Homography Solver (\texttt{homography.py}) \\
\midrule
\textbf{Obj 6: System Assembly} & - End-to-end CPU latency profiling & - Validates real-time execution ($34.4$ FPS) & - Multi-threaded Runtime Engine \\
and Real-Time Execution & - Multi-threaded stream testing & - Confirms XML action dispatching & - Action Command Mapper GUI \\
\bottomrule
\end{tabular}%
\end{table}

\subsection{Qualitative Data Gathering from User Interactivity Studies}
\label{subsec:req_qualitative_data}
The user requirements acquisition process has been performed considering the input provided by 15 test users having different levels of typing proficiency. It was revealed that the utilization of static layouts results in decreased efficiency while working with some desktop applications. Users needed a personalized macro pad, which means that the need for the Layout Designer (\texttt{designer/designer\_app.py}) is evident.

\subsection{Empirical Kinematic Data Translation to Landmark Requirements}
\label{subsec:req_kinematic_translation}
The skeletal data recorded across 15 participants confirmed that finger contact events are accompanied by sharp deceleration spikes in vertical fingertip velocity. Consequently, the feature engineering requirements specify unitless hand-length scale normalization ($L_{\text{hand}}$) without low-pass filtering, directly propagating instantaneous kinematic features into the neural classifier.

\section{Unified Modeling Language (UML) System Analysis}
\label{sec:req_uml_analysis}

\subsection{Use Case Analysis and Formal Specifications}
\label{subsec:req_use_cases}
The system architecture contains two human actors: Layout Designer User (operates with Application 1) and Runtime Operator User (operates with Application 2). Auxiliary actors of the system include the Monocular RGB Camera and the Operating System Shell.

Figure~\ref{fig:use_case_diagram} shows the Use Case Diagram covering system boundaries, actors, and use case interactions.

\begin{figure}[htbp]
\centering

\includegraphics[width=0.85\textwidth]{figures/chapter04_fig02.png}

\caption{System Use Case Diagram covering Application 1 and Application 2 workflows.}
\label{fig:use_case_diagram}
\end{figure}

Detailed formal Use Case Specifications for core system operations are provided in Tables~\ref{tab:uc_spec_01}, \ref{tab:uc_spec_02}, \ref{tab:uc_spec_03}, \ref{tab:uc_spec_04}, \ref{tab:uc_spec_05}, \ref{tab:uc_spec_06}, \ref{tab:uc_spec_07}, and \ref{tab:uc_spec_08}.

\begin{table}[htbp]
\centering
\caption{Formal Use Case Specification: UC-01 Design Custom Key Layout.}
\label{tab:uc_spec_01}
\begin{tabular}{l p{11cm}}
\toprule
\textbf{Use Case Attribute} & \textbf{Specification Detail} \\
\midrule
\textbf{Use Case ID and Title} & \textbf{UC-01: Creation of Custom Key Layout} \\
\textbf{Primary Actor} & Layout Designer User \\
\textbf{Description} & Creation of custom key layout by user is done using placement of keys on PySide6 A4 GUI canvas by providing location of buttons and setting dimensions of keys. \\
\textbf{Preconditions} & Running of Application 1 (\texttt{designer/designer\_app.py}) on host desktop computer. \\
\textbf{Postconditions} & Creation of layout graphical elements and button metadata in GUI memory. \\
\textbf{Main Success Scenario} & 1. User runs Designer GUI and marks boundaries of keys on grid.\\
 & 2. User provides key properties (Button ID, default values, command bindings).\\
 & 3. System places AprilTags at corners of layout automatically.\\
\textbf{Alternative Flows} & 1a. Boundaries of keys are overlapping: System warns user about it and prompts correction. \\
\bottomrule
\end{tabular}%
\end{table}

\begin{table}[htbp]
\centering
\caption{Formal Use Case Specification: UC-02 Export Synchronized PDF/XML Artifacts.}
\label{tab:uc_spec_02}
\begin{tabular}{l p{11cm}}
\toprule
\textbf{Use Case Attribute} & \textbf{Specification Detail} \\
\midrule
\textbf{Use Case ID and Title} & \textbf{UC-02: Export Synchronized PDF and XML Artifacts} \\
\textbf{Primary Actor} & Layout Designer User \\
\textbf{Description} & System generates a vector PDF file with anchors of AprilTags to be printed and generates an XML file with the layout and bounding boxes. \\
\textbf{Preconditions} & Good key layout design without any overlapping of boundaries in UC-01. \\
\textbf{Postconditions} & Generated PDF and XML files will be saved on the disk directory. \\
\textbf{Main Success Scenario} & 1. User initiates the process of Export from Application 1 File menu.\\
 & 2. System verifies the anchors of AprilTags along margins of the layout.\\
 & 3. System generates vector PDF file to be printed and saves XML file on the disk.\\
\textbf{Alternative Flows} & 2a. If there is no disk write access available: System notifies the user and prompts for alternative path. \\
\bottomrule
\end{tabular}%
\end{table}

\begin{table}[htbp]
\centering
\caption{Formal Use Case Specification: UC-03 Select Camera Input Stream.}
\label{tab:uc_spec_03}
\begin{tabular}{l p{11cm}}
\toprule
\textbf{Use Case Attribute} & \textbf{Specification Detail} \\
\midrule
\textbf{Use Case ID and Title} & \textbf{UC-03: Select Camera Video Input Stream} \\
\textbf{Primary Actor} & Runtime Operator User \\
\textbf{Description} & The user picks up the input stream of the currently connected monocular RGB camera from Application 2 Setup GUI. \\
\textbf{Preconditions} & The camera device is connected to the host desktop computer using the USB interface. \\
\textbf{Postconditions} & The camera capture thread (\texttt{camera\_thread.py}) is launched using the specified device index. \\
\textbf{Main Success Scenario} & 1. The operator starts the Application 2 Setup GUI window.\\
 & 2. The system detects available video capture devices.\\
 & 3. The operator picks up the camera device to be used from the dropdown list.\\
 & 4. The system establishes an OpenCV video stream with live preview of images.\\
\textbf{Alternative Flows} & 2a. No cameras detected: The system reports the hardware problem. \\
\bottomrule
\end{tabular}%
\end{table}

\begin{table}[htbp]
\centering
\caption{Formal Use Case Specification: UC-04 Load Layout XML Schema.}
\label{tab:uc_spec_04}
\begin{tabular}{l p{11cm}}
\toprule
\textbf{Use Case Attribute} & \textbf{Specification Detail} \\
\midrule
\textbf{Use Case ID and Title} & \textbf{UC-04: Loading Layout XML Schema} \\
\textbf{Primary Actor} & Runtime Operator User \\
\textbf{Description} & System parses the XML layout schema file and creates lookup tables for target coordinates. \\
\textbf{Preconditions} & The XML layout schema file generated by Application 1 must be saved on disk. \\
\textbf{Postconditions} & The bounding boxes and button IDs are loaded into shared runtime memory. \\
\textbf{Main Success Scenario} & 1. User selects 'Load Layout XML' option from the Application 2 setup user interface.\\
 & 2. The system parses the XML file to get the bounding boxes and AprilTag anchor locations.\\
 & 3. The system creates the spatial data structures for use in real-time key matching.\\
\textbf{Alternative Flows} & 2a. Invalid XML schema: The system shows error validation message. \\
\bottomrule
\end{tabular}%
\end{table}

\begin{table}[htbp]
\centering
\caption{Formal Use Case Specification: UC-05 Map Buttons to Keystrokes/Commands.}
\label{tab:uc_spec_05}
\begin{tabular}{l p{11cm}}
\toprule
\textbf{Use Case Attribute} & \textbf{Specification Detail} \\
\midrule
\textbf{Use Case ID and Title} & \textbf{UC-05: Map Key Strokes or Shell Commands to Button Identifier} \\
\textbf{Primary Actor} & Runtime Operator User \\
\textbf{Description} & User assigns unique button identifiers with individual keystrokes or system shell commands. \\
\textbf{Preconditions} & XML file structure has been successfully loaded in UC-04. \\
\textbf{Postconditions} & Action map table has been successfully registered by the runtime execution engine. \\
\textbf{Main Success Scenario} & 1. The system assigns button identifiers into the interactive action map table.\\
 & 2. The user assigns button identifiers with keystrokes (e.g., `'A'`, `'Enter'`) or shell scripts.\\
 & 3. The system registers action bindings in the dispatcher queue.\\
\textbf{Alternative Flows} & 2a. Error in shell script syntax: The system flags the error command string. \\
\bottomrule
\end{tabular}%
\end{table}

\begin{table}[htbp]
\centering
\caption{Formal Use Case Specification: UC-06 Save and Load Action Configuration Profiles.}
\label{tab:uc_spec_06}
\begin{tabular}{l p{11cm}}
\toprule
\textbf{Use Case Attribute} & \textbf{Specification Detail} \\
\midrule
\textbf{Use Case ID and Title} & \textbf{UC-06: Save and Load Action Combination Profiles} \\
\textbf{Primary Actor} & Runtime Operator User \\
\textbf{Description} & The user saves the command mappings to disk in a reusable JSON file format and switches between different action combination profiles using the exact same physical printed layout sheet without reprinting the paper. \\
\textbf{Preconditions} & Creation of action mapping table in UC-05. \\
\textbf{Postconditions} & Saving or loading of a profile JSON file either on the disk or active mapping table. \\
\textbf{Main Success Scenario} & 1. User clicks on 'Save Profile' button in Application 2 GUI.\\
 & 2. System creates a JSON file from the action mapping table (e.g., \texttt{Work\_Mode.json}).\\
 & 3. User switches between different profile files (e.g., \texttt{DAW\_Mode.json}, \texttt{Gaming\_Mode.json}) using the same physical printed layout sheet without reprinting.\\
\textbf{Alternative Flows} & 3a. Different profile file format from the existing one: System performs automatic migration. \\
\bottomrule
\end{tabular}%
\end{table}

\begin{table}[htbp]
\centering
\caption{Formal Use Case Specification: UC-07 AprilTag Tracking and Planar Homography Estimation.}
\label{tab:uc_spec_07}
\begin{tabular}{l p{11cm}}
\toprule
\textbf{Use Case Attribute} & \textbf{Specification Detail} \\
\midrule
\textbf{Use Case ID and Title} & \textbf{UC-07: Track AprilTag Anchors and Calculate Planar Homography Matrix $H$} \\
\textbf{Primary Actor} & Monocular RGB Camera (Hardware) / Background Runtime Engine \\
\textbf{Description} & Detecting positions of AprilTag anchors in the captured camera image and calculating the $3 \times 3$ planar homography matrix $H$ through SVD. \\
\textbf{Preconditions} & Camera feed is active; printed paper layout is visible to the camera. \\
\textbf{Postconditions} & The planar homography matrix $H \in \mathbb{R}^{3 \times 3}$ is updated in shared runtime memory. \\
\textbf{Main Success Scenario} & 1. Video stream captured by the camera thread is provided to the AprilTag thread.\\
 & 2. Positions of AprilTag anchors are detected by AprilTag detector with sub-pixel accuracy.\\
 & 3. The detected image anchors are matched against the anchors defined in the XML file.\\
 & 4. SVD is computed solving homography equations ($\mathbf{A} \mathbf{h} = \mathbf{0}$) and homography matrix $H$ is updated.\\
\textbf{Alternative Flows} & 3a. Tags are covered by hands (fewer than 4 tags detected): Previously computed valid homography matrix is used. \\
\bottomrule
\end{tabular}%
\end{table}

\begin{table}[htbp]
\centering
\caption{Formal Use Case Specification: UC-08 Hand Landmark Scale Normalization and LSTM Touch Classification.}
\label{tab:uc_spec_08}
\begin{tabular}{l p{11cm}}
\toprule
\textbf{Use Case Attribute} & \textbf{Specification Detail} \\
\midrule
\textbf{Use Case ID and Title} & \textbf{UC-08: Normalized Keypoint Coordinates, Classification by PyTorch LSTM, and Dispatch Action} \\
\textbf{Primary Actor} & Background Runtime Engine / OS Shell Execution Engine \\
\textbf{Description} & System detects 21 MediaPipe hand keypoints, normalizes coordinates according to hand length $L_{\text{hand}}$, uses 5-frame sliding window and PyTorch LSTM classification, translates active fingertip coordinate to XML key coordinate space using homography matrix $H$, and executes an action based on keypress or mapped command. \\
\textbf{Preconditions} & Layout XML is loaded; PyTorch model (\texttt{best\_finger\_touch\_lstm.pth}) is loaded in memory. \\
\textbf{Postconditions} & Keystroke or OS shell command is performed; active touch status is reported in GUI. \\
\textbf{Main Success Scenario} & 1. MediaPipe detects 21 2D joints' keypoints from RGB image from the camera.\\
 & 2. System normalizes keypoints according to wrist and middle MCP hand length $L_{\text{hand}}$.\\
 & 3. System creates 5-frame sliding window and passes matrix to PyTorch LSTM.\\
 & 4. PyTorch LSTM classifies active finger surface touch event ($P > 0.90$).\\
 & 5. System translates active fingertip pixel $\mathbf{P}_{\text{pixel}}$ to XML coordinate space ($\mathbf{P}_{\text{XML}} = H \cdot \mathbf{P}_{\text{pixel}}$).\\
 & 6. Key lookup retrieves ID of the hit key and executes mapped OS shell command.\\
\textbf{Alternative Flows} & 4a. Motion of hovering kind detected ($P \le 0.90$): System blocks key activation. \\
\bottomrule
\end{tabular}%
\end{table}

\subsection{System Class Diagram and Class Responsibility Breakdown}
\label{subsec:req_class_diagram}
The architecture of the system has been developed under OOP. In this respect, the different functionalities of the system have been separated into different classes, including Layout, Camera, Feature, Deep Learning, Homography, and Command Execution. Figure~\ref{fig:class_diagram} illustrates the System Class Diagram.

\begin{figure}[htbp]
\centering

\includegraphics[width=0.85\textwidth]{figures/chapter04_fig03.png}

\caption{System Class Diagram covering Layout Designer and Runtime Keyboard Engine modules.}
\label{fig:class_diagram}
\end{figure}

The class responsibility collaboration (CRC) for the core classes of the system is described below:
\begin{enumerate}
    \item \textbf{\texttt{LayoutDesignerApp}:} Manages PySide6 window and provides graphics buttons as well as export features. Works in coordination with \texttt{KeyButtonGraphic} and export utilities.
    \item \textbf{\texttt{RuntimeKeyboardEngine}:} Handles vision processing from the back end, works with threads, and coordinates hand tracking and key actions. Works with \texttt{CameraThread}, \texttt{MediaPipePoseExtractor}, and \texttt{KeyActionDispatcher}.
    \item \textbf{\texttt{ScaleNormalizationEngine}:} Shifts joint coordinates relative to the wrist $\mathbf{P}_0$ and computes dimensionless scaling factors based on hand length $L_{\text{hand}}$. Works with \texttt{SlidingWindowBuffer}.
    \item \textbf{\texttt{HomographyMappingEngine}:} Finds AprilTag corner points and computes the $H \in \mathbb{R}^{3 \times 3}$ matrix using SVD techniques. Works with \texttt{AprilTagTracker} and \texttt{KeyActionDispatcher}.
\end{enumerate}

\subsection{System Activity Diagram}
\label{subsec:req_activity_diagram}
Multiple threads are used by the runtime engine to achieve real-time behavior using standard commodity CPUs. Figure~\ref{fig:activity_diagram} presents the System Activity Diagram where various processes run in parallel for video capture, AprilTag localization, hand pose recognition, sliding window LSTM touch recognition, and homography processing.

\begin{figure}[htbp]
\centering

\includegraphics[width=0.85\textwidth]{figures/chapter04_fig04.png}

\caption{System Activity Diagram illustrating parallel vision execution threads.}
\label{fig:activity_diagram}
\end{figure}

\subsection{Sequence Diagrams Mapped to Key Use Cases}
\label{subsec:req_sequence_diagrams}
Sequence diagrams are employed to depict the flow of messages amongst the different software components during operation. Figure~\ref{fig:seq_diagram_designer} shows Sequence Diagram 1 depicting the layout design and synchronized file export for UC-01 and UC-02.

\begin{figure}[htbp]
\centering

\includegraphics[width=0.85\textwidth]{figures/chapter04_fig05.png}

\caption{Sequence Diagram 1: Layout Design and Synchronized Export Sequence (UC-01 and UC-02).}
\label{fig:seq_diagram_designer}
\end{figure}

Figure~\ref{fig:seq_diagram_runtime} shows Sequence Diagram 2 for real-time vision, AprilTag tracking, LSTM inference, and command execution for UC-07 and UC-08.

\begin{figure}[htbp]
\centering

\includegraphics[width=0.85\textwidth]{figures/chapter04_fig06.png}

\caption{Sequence Diagram 2: Real-Time Vision, AprilTag Tracking, PyTorch LSTM Inference, and Action Execution Sequence (UC-07 and UC-08).}
\label{fig:seq_diagram_runtime}
\end{figure}

\subsection{System Statechart Diagram}
\label{subsec:req_statechart_diagram}
The states in the lifecycle of the finger touch interaction are illustrated through the Statechart diagram shown in Figure~\ref{fig:statechart_diagram}.

\begin{figure}[htbp]
\centering

\includegraphics[width=0.85\textwidth]{figures/chapter04_fig07.png}

\caption{System Statechart Diagram governing finger touch lifecycle transitions.}
\label{fig:statechart_diagram}
\end{figure}

\subsection{System Deployment Diagram}
\label{subsec:req_deployment_diagram}
Figure~\ref{fig:deployment_diagram} presents the System Deployment Diagram which illustrates the hardware elements, execution environment, and software elements.

\begin{figure}[htbp]
\centering

\includegraphics[width=0.85\textwidth]{figures/chapter04_fig08.png}

\caption{System Deployment Diagram detailing hardware nodes, execution processes, and software artifacts.}
\label{fig:deployment_diagram}
\end{figure}

\section{Proposed System Architecture Framework}
\label{sec:req_architecture}
The architecture framework for the proposed system is organized as a two-application software suite. The data flows present in the architecture framework between Application 1 (Layout Designer Suite) and Application 2 (Runtime Virtual Keyboard Engine and Command Mapper) are shown in Figure~\ref{fig:system_architecture_framework}.

\begin{figure}[htbp]
\centering

\includegraphics[width=0.85\textwidth]{figures/chapter04_fig09.png}

\caption{Proposed System Architecture Framework Diagram highlighting Application 1 and Application 2 synchronized workflows.}
\label{fig:system_architecture_framework}
\end{figure}

As shown in Figure~\ref{fig:system_architecture_framework}, the design separation keeps layout authoring distinct from live camera processing:
\begin{itemize}
    \item \textbf{Application 1 (Layout Designer Suite):} Desktop GUI application that generates a printable vector PDF layout with AprilTag border markers on plain paper along with XML files of layout coordinate definitions.
    \item \textbf{Application 2 (Runtime Keyboard Engine):} Contains a Setup and Command Mapper GUI (Component A) for defining key-to-command mappings and choosing the camera stream, and a Background Vision Watch Engine (Component B) that detects AprilTag homography $H$, extracts MediaPipe hand landmarks, performs scale normalization ($L_{\text{hand}}$), evaluates a 12~FPS 5-frame sliding window using PyTorch LSTM on CPU, transforms coordinates ($\mathbf{P}_{\text{XML}} = H \cdot \mathbf{P}_{\text{pixel}}$), and dispatches key commands.
\end{itemize}

\subsection{Architectural Design Patterns Applied}
\label{subsec:req_design_patterns}
For the code to remain modular and maintainable, the following design patterns are implemented in the software architecture:
\begin{enumerate}
    \item \textbf{Model-View-Controller (MVC):} Used in the Layout Designer (\texttt{designer\_app.py}) where the visual canvas (View) is separated from XML coordinate data (Model) and user mouse actions (Controller).
    \item \textbf{Producer-Consumer Thread Queue:} Used in Application 2, where video frames are captured by \texttt{camera\_thread.py} and queued in a thread-safe manner for the model manager (\texttt{model\_manager.py}) to avoid frame loss.
    \item \textbf{Strategy Pattern for Action Dispatching:} Separates regular keyboard typing from shell script execution into modular command handlers.
    \item \textbf{Singleton Model Manager:} Ensures the pre-trained PyTorch touch model (\texttt{best\_finger\_touch\_lstm.pth}) is loaded into memory only once during startup.
\end{enumerate}

\section{Functional and Non-Functional Requirements}
\label{sec:req_fn_nfn}
There are two categories of system requirements: functional requirements (FR) and non-functional requirements (NFR). The functional requirements specify what basic functions can be achieved by the system and are outlined in Table~\ref{tab:functional_requirements} below.

\subsection{Functional Requirements (FR)}
\label{subsec:req_functional}

\begin{table}[htbp]
\centering
\caption{Detailed Functional Requirements (FR) specification matrix.}
\label{tab:functional_requirements}
\begin{tabular}{l l p{8cm}}
\toprule
\textbf{Req ID} & \textbf{Requirement Title} & \textbf{Detailed Functional Requirement Description} \\
\midrule
\textbf{FR-01} & Interactive Layout Design & The PySide6 GUI interface shall enable the user to make and manipulate key buttons like create, resize, drag and drop, move, and delete on an A4 layout grid ($210 \times 297\text{ mm}$). \\
\textbf{FR-02} & AprilTag Anchor Location & The Designer application shall automatically calculate and place the AprilTag visual tags at the four corners of the border layout. \\
\textbf{FR-03} & Vector PDF Layout Renderer & Application 1 shall render and export the vector PDF layout that can be physically printed on paper with AprilTags. \\
\textbf{FR-04} & XML Schema Generator & Application 1 shall export the structural XML file containing key IDs, bounding boxes $(x_{\min}, y_{\min}, x_{\max}, y_{\max})$, default values, and tag locations. \\
\textbf{FR-05} & Regular RGB Camera Stream Selection & Application 2 Setup GUI shall enable the detection and selection of a single regular monocular RGB camera stream across regular and low resolutions, without specialized hardware. \\
\textbf{FR-06} & XML Specifier Loader & Application 2 shall process the exported XML layout file to generate the target coordinate lookup table. \\
\textbf{FR-07} & Key Press Action Mapper & The Setup GUI shall enable the user to map key IDs with single key presses such as \texttt{'A'}, \texttt{'Space'}, \texttt{'Enter'}. \\
\textbf{FR-08} & System Command Action Mapper & The Setup GUI shall enable the user to map key IDs with system shell commands or external Python scripts. \\
\textbf{FR-09} & Action Profile Multiplexing & The system shall enable the user to store, load, and switch between multiple distinct action profile combinations for the exact same physical printed layout sheet without reprinting. \\
\textbf{FR-10} & Multi-Threaded Video Frame Capture & Runtime engine must capture OpenCV video frames in a dedicated background thread to avoid dropping frames. \\
\textbf{FR-11} & AprilTag Sub-Pixel Corner Localization & AprilTag tracking engine must localize quad tags and determine sub-pixel corner coordinates when the camera is tilted. \\
\textbf{FR-12} & SVD Homography Matrix Calculation & Engine must compute the $3 \times 3$ Planar Homography matrix $H$ with the use of SVD on AprilTag correspondences. \\
\textbf{FR-13} & MediaPipe 21 Joint Pose Estimation & Engine must determine 2D pixel coordinates of 21 hand joint keypoints (\texttt{num\_hands=1}) for each video frame. \\
\textbf{FR-14} & Unitless Hand Scale Normalization & System must normalize hand joint coordinates with respect to the origin of the wrist joint $\mathbf{P}_0$ and scale them to hand length $L_{\text{hand}} = \|\mathbf{P}_9 - \mathbf{P}_0\|_2$. \\
\textbf{FR-15} & Direct Feature Propagation & Feature extraction pipeline must provide scale-normalized hand joint coordinates directly to the classifier without temporal low-pass filtering. \\
\textbf{FR-16} & 5-Frame Sliding Window Construction & System must build 84-dimensional feature vectors from 5 frames with 2-frame overlap (stride of 3 frames) at 12~FPS for the active single hand. \\
\textbf{FR-17} & PyTorch LSTM Touch Inference & Lightweight PyTorch model (\texttt{best\_finger\_touch\_lstm.pth}) must classify finger paper touch contact probability ($P > 0.90$) across all 5 digits on standard CPU. \\
\textbf{FR-18} & Planar Coordinate Transformation & Engine must transform active hand fingertip coordinates to XML layout space using $H$ ($\mathbf{P}_{\text{XML}} = H \cdot \mathbf{P}_{\text{pixel}}$) and execute action. \\
\bottomrule
\end{tabular}%
\end{table}

\subsection{Detailed Functional Requirements Subsection Breakdown}
\label{subsec:req_fr_details}
There are 18 functional requirements grouped according to four main subsystems:
\begin{enumerate}
    \item \textbf{Layout Designer Domain (FR-01 to FR-04):} Implemented in Application 1 (\texttt{designer/designer\_app.py}). Ensures reliable interaction with button creation, bounding box calculation, AprilTag anchor point positioning, vector PDF rendering, and XML layout file generation for any custom keyboard layout.
    \item \textbf{Setup and Command Mapper Domain (FR-05 to FR-09):} Implemented in Application 2 Component A. Provides selection of cameras from standard and low resolutions, XML layout loading, assignments of custom keystrokes and shell commands, and saving/loading profile files which allows switching between action combinations on the exact same printed paper sheet without reprinting.
    \item \textbf{Vision Perception and Scale Normalization Domain (FR-10 to FR-16):} Implemented in Application 2 Component B. Includes multi-threaded camera capture (\texttt{camera\_thread.py}), AprilTag localization (\texttt{aprilTag/estimater.py}), MediaPipe single-hand joint tracking (\texttt{num\_hands=1}), hand length scale normalization ($L_{\text{hand}}$) with direct feature propagation, and 12~FPS 5-frame sliding window assembly (\texttt{create\_windows.py}).
    \item \textbf{Touch Classification and Command Execution Domain (FR-17 to FR-18):} Implemented in Application 2 Component B. Covers CPU-based PyTorch LSTM touch classification (\texttt{best\_finger\_touch\_lstm.pth}) across all 5 fingers of the active hand, planar homography coordinate transformation ($\mathbf{P}_{\text{XML}} = H \cdot \mathbf{P}_{\text{pixel}}$), XML key detection, and keystroke/shell command execution.
\end{enumerate}

\subsection{Non-Functional Requirements (NFR)}
\label{subsec:req_non_functional}
The non-functional requirements are comprised of quality requirements, constraints, and performance targets listed in Table~\ref{tab:non_functional_requirements} below.

\begin{table}[htbp]
\centering
\caption{Detailed Non-Functional Requirements (NFR) specification matrix.}
\label{tab:non_functional_requirements}
\begin{tabular}{l l p{8cm}}
\toprule
\textbf{Req ID} & \textbf{Quality Attribute} & \textbf{Detailed Non-Functional Performance Constraint} \\
\midrule
\textbf{NFR-01} & Execution Latency & Time to process a frame should not be more than $83.3\text{ ms}$ ($\ge 12\text{ FPS}$) executed in near real-time mode on commodity CPUs. \\
\textbf{NFR-02} & Touch Classification Accuracy & F1-score of PyTorch LSTM classifier should be at least $0.95$ across all five active fingers. \\
\textbf{NFR-03} & Sub-Pixel Tracking Precision & The error of AprilTag corner localization should be at most $0.05\text{ pixels}$ for perspective inclination angles up to $75^\circ$. \\
\textbf{NFR-04} & Scale Distance Invariance & The system shall guarantee consistent touch detection precision across camera operating distances ranging from 35 cm to 65 cm. \\
\textbf{NFR-05} & Memory Footprint & The size of PyTorch LSTM model should not be greater than $1\text{ MB}$; runtime memory footprint should be less than $500\text{ MB}$. \\
\textbf{NFR-06} & CPU Compute Footprint & The runtime engine should allow execution of near real-time vision loops ($\ge 12\text{ FPS}$) on commodity CPUs without using a GPU. \\
\textbf{NFR-07} & Portability & The software application suite should be portable across Linux (Ubuntu/Debian) and Windows operating systems. \\
\textbf{NFR-08} & Data Privacy and GDPR & The system shall process video frames locally; raw frames from the camera are permanently discarded after extracting anonymized joint coordinates. \\
\textbf{NFR-09} & Illumination Robustness & The vision pipeline should operate properly across indoor illumination levels ranging from $150\text{ lux}$ (dim illumination) to $2,500+\text{ lux}$ (sunlight). \\
\textbf{NFR-10} & Usability & Interactive layout design in Application 1 should allow non-technical users to create custom keypads in less than 5 minutes. \\
\bottomrule
\end{tabular}%
\end{table}

\subsection{Detailed Non-Functional Quality Constraints}
\label{subsec:req_nfr_details}
The 10 non-functional quality constraints impose fundamental performance requirements:
\begin{itemize}
    \item \textbf{Frame Processing Speed and Real-Time Execution (NFR-01, NFR-06):} The system provides an end-to-end frame processing time of under $30\text{ ms}$ (well within the 12~FPS budget of $83.3\text{ ms}$) on standard commodity CPUs without requiring a dedicated GPU.
    \item \textbf{Accuracy of Touch Classification (NFR-02):} The PyTorch LSTM neural network, trained on 5-frame windows of scale-normalized features, achieves an F1-score of $0.963$ across all five fingers.
    \item \textbf{Tracking Accuracy and Camera Tilt Support (NFR-03, NFR-04):} AprilTag corner tracking reaches sub-pixel accuracy of $0.038\text{ px}$, maintaining homography stability under camera tilt angles up to $75^\circ$.
    \item \textbf{Protection Against Data Leakage (NFR-08):} The system processes video streams locally in memory, keeping only anonymized 2D joint coordinates ($(x_i^{\text{norm}}, y_i^{\text{norm}})$) and deleting raw camera frames permanently.
\end{itemize}

\section{Chapter Summary}
\label{sec:req_summary}
In this chapter, the System Requirement Specification (SRS) and analysis using Unified Modeling Language (UML) have been presented for the customizable paper virtual keyboard system.

Section~\ref{sec:req_stakeholder} described primary, secondary, and system stakeholders along with a Power-Interest Matrix and design trade-off resolutions. Section~\ref{sec:req_operationalization} explained how research findings and data collection were transformed into system requirements. Section~\ref{sec:req_uml_analysis} provided UML diagrams including Use Case diagrams with specifications for 8 core use cases, Class diagrams with CRC descriptions, Activity diagrams, Sequence diagrams, Statechart diagrams, and Deployment diagrams. Section~\ref{sec:req_architecture} described the System Architecture Framework Diagram and architectural design patterns. Section~\ref{sec:req_fn_nfn} specified 18 functional requirements and 10 non-functional requirements with domain breakdowns.

The requirements and architectural models established in this chapter provide direct guidelines for the software implementation, algorithm design, and evaluation presented in the subsequent chapters.
